% This file is part of the modular source for Stratified Relational Workspace Theory.
% !TeX root = ../Stratified_Relational_Workspace_Theory.tex
\theorypart{XXVI}{Discussion, Limitations, and Philosophical Consequences}

\topic{94}{Decomposition Relativity}

Part~XVII turns decomposition relativity from a caveat into a formal part of the
certificate. The central lesson is negative and positive at once. No claim of
relationality is invariant under arbitrary redescription. An indiscrete merge
makes the whole state local, while a secret-sharing split can make a centralized
encoder look irreducibly distributed. Yet exact structural content can remain
invariant under faithful coarse-graining, and gate margins can remain stable
through a declared family of operationally meaningful decompositions.

The resulting position is neither that decomposition is arbitrary nor that one
partition is uniquely privileged. Decompositions are constrained by access,
causal interfaces, intervention separability, temporal continuity, and explicit
shared resources. A certificate established only at one interface is a valid
but limited claim. A stronger claim reports the full envelope
\[
\left(
\widetilde m_q,
\underline m_q,
\overline m_q,
\Delta_q^{\mathrm{dec}}
\right)
\]
for each required gate, where \(\tilde m_q\) is the reference margin and the
remaining terms summarize its behavior across \(\mathfrak D\).

This does not eliminate judgment. It makes the judgment inspectable. Analysts
can disagree about the admissible family and identify exactly which merge,
split, hidden resource, or intervention assumption changes the conclusion.
That is preferable to treating a convenient node diagram as if it were part of
nature itself.

\subtopic{94.1 Specification relativity and epistemic restraint}

Decomposition is only one coordinate of the scientific claim. The downstream
target, loss, temporal horizon, intervention family, null panel, and threshold
rule also determine what is being certified. A workspace may be operationally
real for a month-long theorem-development trajectory yet irrelevant to a
second-scale motor target. Those are distinct specifications, not competing
answers to one underspecified question.

The framework therefore distinguishes truth under a declared specification
from robustness across an admissible family. It also distinguishes robustness
inside one fitted realization from identification across all realizations
compatible with the available experiments. A positive point estimate at one
favorable analysis is the weakest admissible conclusion. A stronger claim
requires a positive lower envelope across \(\mathfrak S\), an identified sign
across the compatible realization set, and a finite-sample confidence bound
that remains positive after uniform error control.

This hierarchy is intentionally conservative. It does not demand invariance
under every imaginable redescription or target. It demands that the scientific
universe over which robustness is claimed be declared before certification and
that ambiguity outside the reach of the experiment be reported as partial or
non-identification rather than converted into certainty.

\topic{95}{Realism Without Reification}

The framework is realist about relational organization but does not posit a free-standing substance
between the components. A workspace exists only through a realization, and every change in the
workspace must be implemented by a change in participants, artifacts, channels, or their states. The
quotient
\[
\Workspace=[\mathfrak R]_{\simeq}
\]
is therefore a mathematical account of identity across realizations, not the name of an additional
physical object.

This position allows two claims that are sometimes treated as incompatible. First, the workspace is
physically dependent: no realization, no workspace. Second, it is not identical with any particular
carrier when several distinct realizations preserve the same structural and causal organization. The same
logic is familiar whenever an object is identified by invariants rather than by one material instance. A
melody is not a second substance beside its performances, yet it is not identical with one performance.
A proof is not separate from every inscription that carries it, yet its identity can survive changes of
notation and medium. The analogy is limited, but it clarifies the intended form of realism.

The theory is therefore compatible with physicalism and multiple realization. It rejects both the view
that a workspace must be wholly inside one component and the stronger view that the workspace is an
independent entity floating outside its carriers. What is real is the organized pattern of compatible
representation, transfer, causal use, persistence, and reconstruction.

\topic{96}{Location and the Sense of ``Between''}

Asking where a relational workspace is located can invite a category mistake. The local states have
ordinary physical locations. A notebook is on a desk, a model state is implemented on hardware, a
message is transmitted through a channel, and a person's activity is realized in a biological organism.
The workspace as such is not located at an additional point. Its location is distributed in the modest
sense that the relations constituting it span several located carriers.

The word ``between'' is best understood counterfactually rather than spatially. Suppose local marginal
states are held fixed while transfer organization is shuffled, delayed, or replaced by a matched null. If
the joint predictive or causal capacity changes, then an operative property resided in the organization of
the relations rather than in any one local state. Conversely, if a single component or matched storage
surrogate reproduces the complete capacity, the case for a relational workspace fails. The relevant
question is therefore not whether an invisible object occupies the interval between participants. It is
whether the coupling carries indispensable state and supports capacities that disappear when that
coupling is altered.

This account also explains why the workspace can be partially present in each component without being
fully stored in any of them. Each local chart may encode a view, role, or constraint. The global object is
the compatible and causally operative organization of those views. In exact settings this is a global
section. In approximate settings it is a stable low-energy subspace together with the dynamic and causal
conditions developed above.

\topic{97}{Continuity, Contribution, and Shared Ownership}

Relational persistence does not require an unchanging set of participants. A workspace can survive
replacement when the new realization preserves the relevant transfer structure, operation algebra,
causal response, reconstructive capacity, and non-collapse margin. This is continuity through organized
transport, not continuity through an immutable token.

The same point applies to artifacts. A handoff document, theorem register, code repository, or experiment
log is not merely a passive record when it constrains future reasoning and permits reconstruction after
interruption. Such artifacts can become constitutive parts of a workspace realization. Their contribution
is causal and testable: remove or corrupt them, then measure what can no longer be reconstructed or done.

Distributedness does not imply equal contribution. One participant may supply long-horizon goals, another
formal manipulation, and an artifact precise persistent memory. Nor does a system-level description erase
individual authorship or responsibility. The theory describes how capacities are organized; it does not
settle legal, ethical, or scholarly questions of credit. Those questions require a separate account of who
introduced which ideas, who exercised control, who bore risk, and which contributions were replaceable.

The phrase ``shared ownership'' should therefore be used carefully. No individual component need own a
complete copy of the workspace state, yet individuals can still own artifacts, originate arguments, and
retain responsibility for decisions. Relational identity and ordinary ownership are different relations.

\topic{98}{Agency, Consciousness, and Moral Non-Implications}

The agency certificate developed here is functional and domain relative. It concerns persistent goals,
interventional value estimates, feasible action control, damage detection, repair, and bounded maintenance
within a declared environment. Passing those tests would justify a claim of workspace-level agency in that
domain. It would not establish unrestricted autonomy.

No theorem in the framework identifies agency with phenomenal consciousness. A system may maintain a
goal, evaluate counterfactual actions, and repair its own organization without there being anything it is
like to be that system. The converse also remains open: failure to satisfy this engineering certificate
would not prove the absence of consciousness. The framework is designed to prevent functional evidence
from silently becoming a metaphysical conclusion.

Moral status is likewise not inferred. Neither coherence, causal efficacy, relational irreducibility, nor
self-maintenance is by itself a criterion of moral patienthood. Future arguments connecting these properties
to consciousness or welfare would need independent premises and evidence. Until then, the framework can
support careful descriptions of organization and agency without treating a collaborative system as a
person or dismissing the moral standing of the humans who participate in it.

\topic{99}{Limitations}

The theory is deliberately broad, but its conclusions remain bounded by several limitations.

\begin{enumerate}
\item \textbf{Admissible-decomposition burden.} Part~XVII formalizes faithful coarse-graining and
family-robust certification, but the choice of primitive carriers and admissible alternatives remains an
empirical and methodological judgment. A narrow family can still overstate robustness.

\item \textbf{Admissible-specification burden.} Part~XVIII makes the target, loss, horizon,
intervention family, null panel, and thresholds explicit. The admissible family \(\mathfrak S\) still
requires scientific judgment. A family chosen after seeing the results can conceal fragility even when
every reported analysis is technically correct.

\item \textbf{Limited observability and partial identification.} Many practical systems expose messages
and outputs rather than complete local states. Transfer maps, predictive states, reconstruction operators,
and even the certificate sign may therefore be only partially identified. A point claim is unjustified
when both passing and failing realizations remain compatible with the declared experiment.

\item \textbf{Linear and finite-dimensional emphasis.} The clearest structural theorems use finite graphs
and linear Hilbert-space models. Nonlinear, stochastic, symbolic, and category-valued extensions require
additional work.

\item \textbf{Gauge and model ambiguity.} Even with abundant data, local coordinates need not be
identifiable. Learned transfers are often meaningful only up to a gauge class and relative to a selected
model family.

\item \textbf{Causal identification burden.} Workspace interventions must be feasible and matched.
Observational dependence, activation, or predictive accuracy cannot substitute for interventions when a
causal claim is made.

\item \textbf{Certification cost.} The complete certificate requires coalition tests, null models,
replacement studies, longitudinal interventions, admissible specification sweeps, and identified-set or
uniform uncertainty calculations. Exhaustive certification may be impractical in large systems, and
adaptive exploration can invalidate nominal error guarantees unless selection and certification are
separated or valid selective procedures are used.

\item \textbf{Conditional mathematical results.} The theorems establish what follows from stated
assumptions. They do not show that naturally occurring collaborations satisfy those assumptions, and
several constants may be too conservative for practical estimation.

\item \textbf{Open empirical status.} The motivating human--AI collaboration has not been subjected to
the complete frozen certificate. It remains a plausible case study and a source of hypotheses, not a
demonstrated instance of relational agency.

\item \textbf{Open priority boundary.} The literature comparison is broad but cannot prove absolute
firstness. Closely related work may exist under other terminology or appear after the review date.
\end{enumerate}

These limitations do not make the framework vacuous. They specify the conditions under which its claims
can fail and identify the measurements needed to strengthen them.

\topic{100}{Open Problems}

Several directions are mathematically and empirically immediate.

\begin{enumerate}
\item Develop nonlinear and stochastic analogues of the transfer Laplacian, including kernel, manifold,
and category-valued realizations.

\item Establish sharp uniform finite-sample estimators for workspace dimension, coalition risk,
matched-null gaps, replacement identity, and specification envelopes under adaptive data collection.

\item Characterize the identified set when local states are latent and only messages, artifacts, or
actions are observed. Determine the minimum intervention family needed to identify each certificate gate
and the complete conjunctive certificate.

\item Extend gluing and persistence to changing graph topology, asynchronous communication, and
delayed or lossy transfers.

\item Extend the joint decomposition-and-specification theory to latent shared resources, overlapping
component systems, continuous coarse-graining, target families, and data-driven discovery of admissible
interface and specification families without invalidating certification.

\item Study the interaction between workspace formation and strategic participants whose local objectives
need not agree.

\item Build benchmark suites in which structural, causal, persistence, and repair quantities are known
exactly but remain nontrivial to infer.

\item Run preregistered human--AI and multi-agent experiments that estimate the complete certificate
rather than isolated performance gains.

\item Compare certificate-directed self-maintenance with established notions of organizational resilience,
institutional memory, and collective agency.

\item Clarify which additional evidence, if any, could connect functional workspace predicates to theories
of consciousness without assuming the connection in advance.
\end{enumerate}

\topic{101}{Conclusion}

A stratified relational workspace is a temporally persistent, causally efficacious, jointly reconstructible
structure represented only partially in several heterogeneous components. Its identity is determined by
invariants preserved under admissible transfer, and its operative capacity is relational only when no
individual component or matched local-storage null reproduces the relevant function.

The theory gives this definition several independent layers. Exact shared content appears as a global
compatible section and, in the invertible case, as a holonomy-fixed core. Approximate shared content
appears as a low-energy spectral subspace of the transfer Laplacian. Predictive states and feasible
interventions determine whether that structure matters for a declared future. Coalition tests and matched
nulls determine whether the capacity is genuinely distributed. Reconstruction frames, temporal transport,
and identity profiles determine whether it survives loss and replacement. Descent determines whether
local workspaces compose globally. Provenance tests distinguish learned organization from scaffolded
answers. Goal, value, damage, and repair interventions determine whether the workspace supports
domain-relative agency and self-maintenance.

No one layer can stand in for the others. Coherence can describe an archive. Performance can describe a
centralized controller. Persistence can be trivial if all content collapses. A task-equivalent replacement
can alter the causal organization. Local compatibility can fail to glue globally. Endogenous-looking
structure can be supplied through hidden leakage, and successful recovery can be open-loop replay rather
than repair. The value of the framework lies in making these alternatives explicit and testable.


Relationality also cannot be detached from decomposition. Faithful
coarse-graining preserves the exact structural core, but coalition counts and
local sufficiency change when components are merged. Arbitrary splitting can
create a formally irreducible code around a centralized encoder. The strongest
claims therefore quantify their margins across an admissible family of
interfaces rather than relying on one favorable diagram.

A complete scientific claim must also survive scrutiny of the specification and the experiment.
The strongest form of the certificate reports a positive gatewise lower envelope across a predeclared
family of decompositions and scientific specifications, shows that the sign of that envelope is identified
across all realizations compatible with the intervention data, and retains a positive margin after uniform
finite-sample uncertainty. When these conditions fail, the correct conclusion is specification-specific,
fragile, partially identified, or unidentified rather than simply positive or negative.

The central identity principle is therefore
\[
\boxed{
\text{workspace identity}
=
\text{persistent invariant causal organization under admissible transport}.
}
\]
This principle is compatible with changing carriers, representations, and participants, but only when
structural, operational, causal, reconstructive, and non-collapse defects remain within declared tolerances.
It provides a formal language for studying distributed collaboration without turning that collaboration
into a mysterious third substance and without inferring consciousness from organization alone.
